% GENERATED FILE. Edit canonical lessons, then run npm run generate:books.
% canonical-source-hash: fnv1a64:7b2b3ccbb61a58a1
% canonical-lessons: BN-C74-dami, BN-C74-sosta, BN-C74-nichu, BN-C74-bhorti, BN-C74-khali

\chapter{Expensive, Cheap, Low, Full, Empty}
\label{ch:bn-expensive-cheap-low-full-empty}
\hlchaptermodality{74}

\begin{chapteropening}
\textbf{By the end of this chapter:} \emph{I can say expensive, cheap, low, full and empty.}

Complete the last lesson of chapter 74: I can say expensive, cheap, low, full and empty.
\end{chapteropening}

\section[\texorpdfstring{\bn{দামি} (dāmi)}{দামি (dāmi)}]{\bn{দামি} (dāmi) — expensive}

\label{lesson:BN-C74-dami}



\begin{quote}
Before the new one: say the Bengali for quickly, then the Bengali for easy.
\end{quote}

\subsection*{You'll want to know: \bn{দামি}}
\textbf{\bn{দামি}} — \emph{dāmi} — \textquotedblleft{}expensive\textquotedblright{}.

\begin{cousinweb}
\textbf{\bn{দামি}} (\emph{dāmi}) is \emph{dām}, \textquotedblleft{}price\textquotedblright{}, with the ending \emph{-i}. CDIAL lists Bengali \emph{dām} under Sanskrit \emph{dramma}, \textquotedblleft{}a coin\textquotedblright{}, which it traces to Greek \emph{drakhmē} through Middle Persian (CDIAL 6622).
\end{cousinweb}

\begin{grammarlens}[title={What you've built}]
One more word for this chapter.
\end{grammarlens}

\subsection*{Your turn}
\begin{itemize}
\raggedright
  \item \emph{Say it:} \emph{dāmi}
  \item \emph{Say it:} \emph{dāmi}, once more
  \item \emph{Say it:} say \emph{dāmi}
  \item \emph{Recall:} say the Bengali for quickly, then the Bengali for easy, then say \emph{dāmi} again
\end{itemize}

\begin{tcolorbox}[breakable,colback=teal!4,colframe=teal!35!black,title={Before you move on}]
What does \bn{দামি} mean? (\textquotedblleft{}Expensive\textquotedblright{}.) Say it once more.
\end{tcolorbox}
\section[\texorpdfstring{\bn{সস্তা} (sôstā)}{সস্তা (sôstā)}]{\bn{সস্তা} (sôstā) — cheap}

\label{lesson:BN-C74-sosta}



\begin{quote}
Before the new one: say the Bengali for easy, then the Bengali for expensive.
\end{quote}

\subsection*{You'll want to know: \bn{সস্তা}}
\textbf{\bn{সস্তা}} — \emph{sôstā} — \textquotedblleft{}cheap\textquotedblright{}.

\begin{cousinweb}
The origin of \textbf{\bn{সস্তা}} (\emph{sôstā}) is not settled. CDIAL has no entry for it, and a search of its word lists by form and by the gloss \textquotedblleft{}cheap\textquotedblright{} turns up nothing of this shape; the same word in Hindi, \emph{sastā}, is untraced there too.
\end{cousinweb}

\begin{grammarlens}[title={What you've built}]
One more word for this chapter.
\end{grammarlens}

\subsection*{Your turn}
\begin{itemize}
\raggedright
  \item \emph{Say it:} \emph{sôstā}
  \item \emph{Say it:} \emph{sôstā}, once more
  \item \emph{Say it:} say \emph{sôstā}
  \item \emph{Recall:} say the Bengali for easy, then the Bengali for expensive, then say \emph{sôstā} again
\end{itemize}

\begin{tcolorbox}[breakable,colback=teal!4,colframe=teal!35!black,title={Before you move on}]
What does \bn{সস্তা} mean? (\textquotedblleft{}Cheap\textquotedblright{}.) Say it once more.
\end{tcolorbox}
\section[\texorpdfstring{\bn{নিচু} (nichu)}{নিচু (nichu)}]{\bn{নিচু} (nichu) — low}

\label{lesson:BN-C74-nichu}



\begin{quote}
Before the new one: say the Bengali for expensive, then the Bengali for cheap.
\end{quote}

\subsection*{You'll want to know: \bn{নিচু}}
\textbf{\bn{নিচু}} — \emph{nichu} — \textquotedblleft{}low\textquotedblright{}.

\begin{cousinweb}
\textbf{\bn{নিচু}} (\emph{nichu}) is a tadbhava, a word inherited from Sanskrit and reshaped by sound change. It comes from \emph{nīca}, \textquotedblleft{}low\textquotedblright{} (CDIAL 7540). CDIAL writes the Bengali form \emph{nica}.
\end{cousinweb}

\begin{grammarlens}[title={What you've built}]
One more word for this chapter.
\end{grammarlens}

\subsection*{Your turn}
\begin{itemize}
\raggedright
  \item \emph{Say it:} \emph{nichu}
  \item \emph{Say it:} \emph{nichu}, once more
  \item \emph{Say it:} say \emph{nichu}
  \item \emph{Recall:} say the Bengali for expensive, then the Bengali for cheap, then say \emph{nichu} again
\end{itemize}

\begin{tcolorbox}[breakable,colback=teal!4,colframe=teal!35!black,title={Before you move on}]
What does \bn{নিচু} mean? (\textquotedblleft{}Low\textquotedblright{}.) Say it once more.
\end{tcolorbox}
\section[\texorpdfstring{\bn{ভর্তি} (bhôrti)}{ভর্তি (bhôrti)}]{\bn{ভর্তি} (bhôrti) — full}

\label{lesson:BN-C74-bhorti}



\begin{quote}
Before the new one: say the Bengali for cheap, then the Bengali for low.
\end{quote}

\subsection*{You'll want to know: \bn{ভর্তি}}
\textbf{\bn{ভর্তি}} — \emph{bhôrti} — \textquotedblleft{}full\textquotedblright{}.

\begin{cousinweb}
\textbf{\bn{ভর্তি}} (\emph{bhôrti}) is built on the stem \emph{bhar-}, \textquotedblleft{}fill\textquotedblright{}: CDIAL lists Bengali \emph{bharā}, \textquotedblleft{}to load\textquotedblright{}, under Sanskrit \emph{bharati}, \textquotedblleft{}bears, fills\textquotedblright{} (CDIAL 9397), but has no entry for \textbf{\bn{ভর্তি}} (\emph{bhôrti}) itself.
\end{cousinweb}

\begin{grammarlens}[title={What you've built}]
One more word for this chapter.
\end{grammarlens}

\subsection*{Your turn}
\begin{itemize}
\raggedright
  \item \emph{Say it:} \emph{bhôrti}
  \item \emph{Say it:} \emph{bhôrti}, once more
  \item \emph{Say it:} say \emph{bhôrti}
  \item \emph{Recall:} say the Bengali for cheap, then the Bengali for low, then say \emph{bhôrti} again
\end{itemize}

\begin{tcolorbox}[breakable,colback=teal!4,colframe=teal!35!black,title={Before you move on}]
What does \bn{ভর্তি} mean? (\textquotedblleft{}Full\textquotedblright{}.) Say it once more.
\end{tcolorbox}
\section[\texorpdfstring{\bn{খালি} (khāli)}{খালি (khāli)}]{\bn{খালি} (khāli) — empty}

\label{lesson:BN-C74-khali}



\begin{quote}
Before the new one: say the Bengali for low, then the Bengali for full.
\end{quote}

\subsection*{You'll want to know: \bn{খালি}}
\textbf{\bn{খালি}} — \emph{khāli} — \textquotedblleft{}empty\textquotedblright{}.

\begin{cousinweb}
\textbf{\bn{খালি}} (\emph{khāli}) is a loan from Arabic \emph{khālī}, \textquotedblleft{}empty\textquotedblright{}, which reached Bengali through Persian. Its Arabic root is \emph{kh-l-w} (Lane).
\end{cousinweb}

\begin{grammarlens}[title={What you've built}]
One more word for this chapter.
\end{grammarlens}

\subsection*{Your turn}
\begin{itemize}
\raggedright
  \item \emph{Say it:} \emph{khāli}
  \item \emph{Say it:} \emph{khāli}, once more
  \item \emph{Say it:} say \emph{khāli}
  \item \emph{Recall:} say the Bengali for low, then the Bengali for full, then say \emph{khāli} again
\end{itemize}

\begin{tcolorbox}[breakable,colback=teal!4,colframe=teal!35!black,title={Before you move on}]
What does \bn{খালি} mean? (\textquotedblleft{}Empty\textquotedblright{}.) Say it once more.
\end{tcolorbox}
